% ============================================================
% CHAPTER 1 — The 0-Code Revolution: Seeing It in Action
% ============================================================
\chapter{The 0-Code Revolution: Seeing It in Action}\index{0-code}\index{Context Engineering}

\section*{What You Will Learn}
By the end of this chapter you will know:
\begin{itemize}
  \item What a 0-code working session looks like in practice
  \item Why context is everything and how it changes the result
  \item What you will build throughout the book
\end{itemize}

\begin{notebox}
If you have not read the Introduction yet, do so now. It explains the ADLC framework, the comparison with the traditional SDLC, and the three core competencies --- Context Engineering, Risk Design, and Confidence Tagging --- that support the entire book.
\end{notebox}

% ---------------------------------------------------------
\section{A 0-Code Working Session: Live}

You already know the theory from the Introduction. Now let us see how it works \textbf{in practice}: what appears on the screen, what you do, and what the AI does.

\subsection{The setup: 30 seconds}

Open VS Code. Open an empty folder. Open Copilot Chat in Agent Mode. Type:

\begin{lstlisting}[language={}]
Create a Python application that converts temperatures between Celsius,
Fahrenheit, and Kelvin. The user chooses the conversion from an interactive
terminal menu. Add a test file using pytest.
\end{lstlisting}

\subsection{What happens in the next 60 seconds}

Copilot in Agent Mode will typically:
\begin{enumerate}
  \item \textbf{Create} the file \texttt{converter.py} with the conversion functions
  \item \textbf{Create} the file \texttt{test\_converter.py} with the unit tests
  \item \textbf{Run} \texttt{pytest} in the integrated terminal
  \item \textbf{Read} the output and verify that all tests pass
  \item \textbf{Report back} with a summary of what it changed
\end{enumerate}

You wrote three lines of instructions. The AI produced roughly $\sim$80 lines of working, tested code.

The important point is this: the code is often \textbf{generic}. It works, but it may not follow a specific structure, may lack robust validation, and may not document the internal logic clearly.

\subsection{The same task, with context}

Now imagine writing a \texttt{\_CONTEXT.md} file first:

\begin{lstlisting}[language=yaml,title={\texttt{\_CONTEXT.md} for Temperature Converter}]
# Project: Temperature Converter

## Purpose
CLI tool for temperature conversion with a focus on scientific precision.

## Constraints
- Kelvin cannot be negative (0 K = absolute zero)
- Always round to 2 decimal places
- Error messages must be in English
- Naming: snake_case everywhere

## Structure
converter/
├── main.py          ← Entry point with menu
├── conversions.py   ← Pure conversion functions
├── validators.py    ← Input validation
└── tests/
    ├── test_conversions.py
    └── test_validators.py
\end{lstlisting}

And then asking:

\begin{lstlisting}[language={}]
Read the _CONTEXT.md and implement the complete application
following all the specifications.
\end{lstlisting}

The result is \textbf{radically different}: separate modules, absolute-zero validation, edge-case tests, messages in English, and a cleaner internal structure.

\begin{quote}
\textbf{Same request. Completely different result.} The variable that changed is the context.
\end{quote}

% ---------------------------------------------------------
\section{Context Is Everything}

This is the most important principle in the entire book.

\textbf{Vague request:}
\begin{lstlisting}[language={}]
Create an API to manage users
\end{lstlisting}

Result: generic code, probably without validation, authentication, or consistent naming.

\textbf{Request with context:}
\begin{lstlisting}[language={}]
Create a REST API with Express.js for user management.
Structure: src/routes/, src/controllers/, src/models/.
Database: PostgreSQL with Prisma ORM.
Endpoints: GET /users, GET /users/:id, POST /users,
PUT /users/:id, DELETE /users/:id.
Input validation with Zod.
Standard JSON responses: { success: boolean, data: T, error?: string }.
Centralized error handling with middleware.
\end{lstlisting}

Result: structured, validated code with repeatable patterns and a foundation suitable for production.

In Chapter~3 you will learn how to write professional \texttt{\_CONTEXT.md} files. In Chapter~4 you will start building your first real project.

\begin{warnbox}
\textbf{A note on expectations.} This book does not require prior programming experience, but it does require \textbf{curiosity} and a willingness to learn new concepts. Starting from Chapter~6 you will meet terms such as API, middleware, database, authentication, and JWT tokens. The AI writes the code, but your role as a \emph{context architect} requires understanding \emph{what} you are asking it to build --- not \emph{how} to type it line by line.
\end{warnbox}

% ---------------------------------------------------------
\section{What You Will Build with This Book}

This book is a \textbf{progression of real projects}, each one more ambitious than the previous.

\subsection*{Project 1 --- Hello World (Chapter 4)}
Your first program generated entirely by AI. You will learn how to formulate requests, inspect the output, and iterate safely.

\subsection*{Project 2 --- CLI Tool (Chapter 5)}
A complete command-line application with CSV/JSON parsing, arguments, and automated tests. This is your first structured \texttt{\_CONTEXT.md}.

\subsection*{Project 3 --- REST API (Chapter 6)}
A web server with Express.js, CRUD endpoints, validation, and Swagger documentation.

\subsection*{Projects 4--7 --- Full-Stack Web Application with OAuth 2.0 (Chapters 7--10)}
\begin{itemize}
  \item \textbf{Project 4} (Ch.~7): Node.js/Express backend with PostgreSQL and Prisma
  \item \textbf{Project 5} (Ch.~8): OAuth~2.0 authentication with JWT
  \item \textbf{Project 6} (Ch.~9): React frontend with protected routes and your first \texttt{SKILL.md}
  \item \textbf{Project 7} (Ch.~10): end-to-end integration and a complete CRUD dashboard
\end{itemize}

\subsection*{Projects 8--9 --- Flutter Mobile App (Chapters 11--12)}
\begin{itemize}
  \item \textbf{Project 8} (Ch.~11): Flutter setup, first mobile app, and Flutter skill file
  \item \textbf{Project 9} (Ch.~12): mobile OAuth login, synchronized CRUD, and deep links
\end{itemize}

\subsection*{Final project --- Full deployment (Chapters 13--15)}
You will bring everything to production: backend on the cloud, frontend on Vercel, managed PostgreSQL, mobile app release, and CI/CD.

% ---------------------------------------------------------
\section{Requirements to Follow This Book}

\subsection*{What you need to know}
\begin{itemize}
  \item How to use a computer (Windows, macOS, or Linux)
  \item The basic idea of how an application works, even just as a user
  \item Basic technical English, because many commands and tools keep their original names
\end{itemize}

\subsection*{What you do \emph{not} need to know}
\begin{itemize}
  \item You do not need to know Python, JavaScript, Dart, or any other language
  \item You do not need prior experience with React, Node.js, Flutter, or PostgreSQL
  \item You do not need to have used the terminal before
  \item You do not need to know Git in advance
\end{itemize}

\subsection*{What you need}
\begin{itemize}
  \item A computer running Windows~10/11, macOS~10.15+, or Linux
  \item A stable internet connection
  \item A free GitHub account
  \item GitHub Copilot or Claude Code access
  \item VS Code installed
\end{itemize}

\begin{costbox}
GitHub Copilot Individual costs about \$10/month. Claude Code operates on a usage basis through the Anthropic API. For most readers, the cost of completing all projects in the book is comparable to an online course or a few hours of technical coaching.
\end{costbox}

% ---------------------------------------------------------
\section{The Promise of This Book}

By the end of this journey you will have:
\begin{enumerate}
  \item \textbf{Built} a complete web application with backend, frontend, and OAuth~2.0 authentication
  \item \textbf{Built} a Flutter mobile app connected to your backend
  \item \textbf{Deployed} everything to production
  \item \textbf{Without manually writing} the code in the traditional way
\end{enumerate}

And above all, you will have acquired a \textbf{method} --- Context Engineering and the ADLC --- that you can reuse for any future application.

\bigskip
\begin{quote}
\itshape ``In the era of 0-code development, human words lose their role as an informal means of communication and take on the weight and consequences of directives in a high-level procedural programming language.''
\end{quote}
